\chapter{Ανάλυση Σημαντικότητας Χαρακτηριστικών}
\label{ch:fi}

Η ανάλυση σημαντικότητας χαρακτηριστικών (feature importance) προσφέρει ερμηνευσιμότητα στα μοντέλα μηχανικής μάθησης, δείχνοντας τους παράγοντες που διαμορφώνουν τις προβλέψεις~\cite{ke2017lgbm, breiman2001rf}. Στο παρόν κεφάλαιο εξετάζουμε πώς η επιλογή της στρατηγικής πρόβλεψης (Teacher-Forcing, Recursive, MIMO, Direct) μετατοπίζει τη βαρύτητα που αποδίδουν τα μοντέλα σε χρονικές υστερήσεις, εξωγενείς μεταβλητές και ημερολογιακά χαρακτηριστικά. Το ερώτημα δεν είναι ποιο υπόδειγμα προβλέπει καλύτερα, αλλά γιατί κάθε στρατηγική αξιοποιεί διαφορετικά κανάλια πληροφόρησης για να ανασυνθέσει τη δυναμική της αγοράς.

% ─────────────────────────────────────────────────────────────────────────────
\section{Μεθοδολογικό Πλαίσιο και Μαθηματική Θεμελίωση}
\label{sec:fi_methodology}
% ─────────────────────────────────────────────────────────────────────────────

Η αποτίμηση της σημαντικότητας των χαρακτηριστικών πραγματοποιείται μέσω τριών διακριτών μαθηματικών προσεγγίσεων, οι οποίες επιλέγονται με βάση την αρχιτεκτονική δομή του εκάστοτε αλγορίθμου, ώστε να αποτυπώνουν τη συνεισφορά κάθε μεταβλητής στη μείωση του σφάλματος παλινδρόμησης.

\paragraph*{1. Κέρδος Πληροφορίας (Gain --- LightGBM, XGBoost).}
Για τους αλγορίθμους ενίσχυσης κλίσης (gradient boosting), η σημαντικότητα ενός
χαρακτηριστικού ισούται με τη συνολική μείωση του σφάλματος που επιφέρουν οι διαχωρισμοί
κόμβων πάνω σε αυτό~\cite{ke2017lgbm}. Έστω $\mathcal{T}$ το σύνολο των εσωτερικών κόμβων
όλων των δένδρων και $\mathcal{T}_f \subseteq \mathcal{T}$ οι κόμβοι που διαχωρίζουν με
βάση το $f$. Η κανονικοποιημένη σημαντικότητα ορίζεται ως:
\begin{equation}
    \mathit{Importance}_{\mathrm{gain}}(f) =
    \frac{\sum_{t \in \mathcal{T}_f} \Delta\mathcal{L}(t)}
         {\sum_{t \in \mathcal{T}} \Delta\mathcal{L}(t)},
    \label{eq:fi_gain}
\end{equation}
όπου $\Delta\mathcal{L}(t)$ είναι η μείωση του μέσου τετραγωνικού σφάλματος (MSE) που
πετυχαίνει ο διαχωρισμός στον κόμβο $t$. Ο παρονομαστής, που αθροίζει σε όλους τους
κόμβους, κανονικοποιεί ώστε $\sum_{f} \mathit{Importance}_{\mathrm{gain}}(f) = 1$.

\paragraph*{2. Μέση Μείωση Ακαθαρσίας (Mean Decrease Impurity / MDI --- Random Forest).}
Στο Random Forest, η σημαντικότητα ισούται με τη μέση, σταθμισμένη ανά πλήθος δειγμάτων,
μείωση ακαθαρσίας που προκαλούν οι διαχωρισμοί πάνω στο $f$~\cite{breiman2001rf}. Για ένα
δάσος $\mathcal{M}$ δένδρων, με $\mathcal{N}_m^{f}$ τους κόμβους του δένδρου $m$ που
διαχωρίζουν με βάση το $f$:
\begin{equation}
    \mathit{Importance}_{\mathrm{MDI}}(f) =
    \frac{1}{|\mathcal{M}|} \sum_{m \in \mathcal{M}}
    \sum_{t \in \mathcal{N}_m^{f}} \frac{N_t}{N}\, \Delta i(t),
    \label{eq:fi_mdi}
\end{equation}
όπου $N_t$ είναι το πλήθος δειγμάτων στον κόμβο $t$ και $N$ το συνολικό. Η μείωση
ακαθαρσίας του κόμβου ορίζεται ως:
\begin{equation}
    \Delta i(t) = i(t) - \frac{N_{t_L}}{N_t}\, i(t_L) - \frac{N_{t_R}}{N_t}\, i(t_R),
    \label{eq:fi_mdi_delta}
\end{equation}
με $i(t)$ τη διακύμανση (variance) των δειγμάτων στον κόμβο $t$ και $t_L, t_R$ τους
θυγατρικούς κόμβους. Η MDI παρουσιάζει γνωστή μεροληψία υπέρ συνεχών μεταβλητών υψηλής
καρδιναλικότητας, που λαμβάνεται υπόψη κατά την ερμηνεία.

\paragraph*{3. Σημαντικότητα μέσω Μετάθεσης (Permutation --- MLP, SVR).}
Στα μη γραμμικά μοντέλα MLP και SVR, που δεν βασίζονται σε δένδρα, οι έννοιες gain και
MDI δεν ορίζονται. Η μέθοδος μετάθεσης μετρά την αύξηση του σφάλματος όταν οι τιμές ενός
χαρακτηριστικού αναδιατάσσονται τυχαία, καταστρέφοντας τη σχέση του με τον στόχο:
\begin{equation}
    \mathit{Importance}_{\mathrm{perm}}(f) =
    \mathrm{MAE}\!\left(\mathcal{D}^{\pi(f)}\right) - \mathrm{MAE}(\mathcal{D}),
    \label{eq:fi_perm}
\end{equation}
όπου $\mathcal{D}$ είναι το σύνολο επικύρωσης και $\mathcal{D}^{\pi(f)}$ το ίδιο σύνολο με
τις τιμές του $f$ τυχαία αναδιατεταγμένες (μετάθεση $\pi$). Λόγω κόστους, η μετάθεση
εφαρμόζεται σε υποδείγματα (800 για το MLP, 150 για το SVR), με αυξημένη στατιστική
διακύμανση στις εκτιμήσεις.

% ─────────────────────────────────────────────────────────────────────────────
\subsection{Πρωτόκολλο Επεξεργασίας και Παραγωγής Αποτελεσμάτων}
\label{subsec:fi_pipeline}
% ─────────────────────────────────────────────────────────────────────────────

Οι χάρτες θερμότητας προκύπτουν από τα πρωτογενή δεδομένα σημαντικότητας των 15 μοντέλων μέσω αυτοματοποιημένης επεξεργασίας. Από κάθε στήλη μοντέλου απομονώνονται τα τρία κορυφαία χαρακτηριστικά και λαμβάνεται η ένωσή τους. Αν προκύψουν πάνω από δέκα, κρατούνται τα δέκα με την καλύτερη μέση κατάταξη σε όλες τις στήλες, ώστε να μην αποκλείονται μεταβλητές κρίσιμες μόνο για συγκεκριμένες στρατηγικές.

Ακολουθούν δύο τύποι χαρτών. Οι χάρτες ανά μοντέλο (Σχ.~\ref{fig:fi_D_tf_price} κ.ε.) δείχνουν κάθε αλγόριθμο ξεχωριστά, επειδή το LightGBM και το XGBoost μετρώνται με gain, το Random Forest με MDI και τα MLP/SVR με μετάθεση, οι τιμές διαφορετικών στηλών δεν συγκρίνονται σε απόλυτη κλίμακα. Γι' αυτό το χρώμα κανονικοποιείται ανά στήλη και οι αναγραφόμενοι αριθμοί δίνουν το ωμό ποσοστό συμμετοχής ($\%\ \mathit{share}$). Οι συγκεντρωτικοί χάρτες ανά στρατηγική (Σχ.~\ref{fig:fi_C_price},~\ref{fig:fi_C_load}) περιλαμβάνουν μόνο τον πρωταθλητή gain κάθε στρατηγικής, καθώς όλες οι στήλες μετρούν την ίδια ποσότητα, η χρωματική κλίμακα είναι ενιαία και η συμπίεση $\mathit{PowerNorm}(\gamma=0{,}45)$ αναδεικνύει τις μικρές τιμές των εξωγενών παραγόντων δίπλα στην κυριαρχία της υστέρησης μίας ώρας στο Teacher-Forcing. Η σύμβαση αυτή δεν επαναλαμβάνεται στις επιμέρους λεζάντες.

% ═════════════════════════════════════════════════════════════════════════════
% (αφαιρέθηκε κενή διπλή επικεφαλίδα ενότητας — η ανάλυση TF/Rec ακολουθεί παρακάτω)
% ═════════════════════════════════════════════════════════════════════════════

% ═════════════════════════════════════════════════════════════════════════════
\section{Teacher-Forcing και Recursive}
\label{sec:fi_tf_rec}
% ═════════════════════════════════════════════════════════════════════════════

Στις δύο μονοβηματικές στρατηγικές το μοντέλο έχει στη διάθεσή του την υστέρηση
μίας ώρας: πραγματική στο Teacher-Forcing, προβλεπόμενη στη Recursive. Η σύγκριση
των δύο απομονώνει την επίδραση της διάδοσης σφάλματος στην κατανομή βαρύτητας.

\begin{figure}[htbp]
  \centering
  \includegraphics[width=0.95\textwidth]{ch4_feature_importance/fig_fi_D_tf_price.png}
  \caption[Σημαντικότητα ανά μοντέλο: Teacher-Forcing, τιμή]{Σημαντικότητα ανά
           μοντέλο: Teacher-Forcing, τιμή DAM.}
  \label{fig:fi_D_tf_price}
\end{figure}

Η κυριαρχία του \textit{Price\_lag1} είναι πραγματική, αλλά το μέγεθός της είναι
χαρακτηριστικό του αλγορίθμου: το φύλλο-προς-φύλλο LightGBM διοχετεύει $88{,}3\%$
στην υστέρηση μίας ώρας, το XGBoost με περιορισμένο βάθος $63{,}3\%$, ενώ το
Random Forest υπό MDI την αραιώνει σε $11{,}2\%$ με τα \textit{Price\_lag24} και
\textit{Price\_lag2} σε παρόμοιο επίπεδο. Η ισχυρή αυτοσυσχέτιση πρώτης τάξης (PACF lag-1 $\approx 1{,}0$,
Κεφάλαιο~\ref{ch:data}) καθιστά την τιμή της προηγούμενης ώρας σχεδόν επαρκή
πρόβλεψη.Οι αλγόριθμοι διαφέρουν στο πόσο επιθετικά την εκμεταλλεύονται, όχι στο
ποιο χαρακτηριστικό μετράει. Η επιλογή είναι στατιστικά αποδοτική όσο η δυναμική
μένει σταθερή και εύθραυστη σε μεταβολή καθεστώτος.

\begin{figure}[htbp]
  \centering
  \includegraphics[width=0.95\textwidth]{ch4_feature_importance/fig_fi_D_rec_price.png}
  \caption[Σημαντικότητα ανά μοντέλο: Recursive, τιμή]{Σημαντικότητα ανά μοντέλο:
           Recursive, τιμή DAM.}
  \label{fig:fi_D_rec_price}
\end{figure}

Στη Recursive η υστέρηση μίας ώρας παύει να είναι παρατήρηση και γίνεται πρόβλεψη
με συσσωρευμένο σφάλμα. Το LightGBM περικόπτει το βάρος της σε $8{,}9\%$ και το
μεταφέρει στα \textit{Price\_lag2} και \textit{Price\_lag24}. Το XGBoost με
προγραμματισμένη δειγματοληψία πηγαίνει παραπέρα, φορτώνοντας \textit{Price\_lag24}
($18{,}5\%$), κυλιόμενο μέσο 24h ($14{,}0\%$) και τις πυκνές υστερήσεις 4-23
(συλλογικά $17{,}1\%$), τις οποίες έμαθε να εμπιστεύεται ακριβώς επειδή η
δειγματοληψία το εξέθεσε σε αυτο-παραγόμενες τιμές κατά την εκπαίδευση. Η
αναδιανομή αυτή είναι η αποτύπωση της διάδοσης σφάλματος στο επίπεδο των
χαρακτηριστικών.

Στο φορτίο η εικόνα είναι αναμενόμενα παρόμοια, με ακόμη ισχυρότερη αυτοσυσχέτιση.
Στο Teacher-Forcing (Σχ.~\ref{fig:fi_D_tf_load}) το \textit{Load\_lag1} συγκεντρώνει
$92{,}4\%$ (LightGBM) και $84{,}0\%$ (XGBoost), υψηλότερα από την τιμή, καθώς η ζήτηση
είναι ομαλότερο σήμα. Το Random Forest υπό MDI κατανέμει ομαλότερα, με τα
\textit{Load\_lag24} ($16{,}5\%$), \textit{Load\_lag2} ($10{,}8\%$) και
\textit{Load\_lag168} ($10{,}6\%$) να αποτυπώνουν τον ημερήσιο και εβδομαδιαίο κύκλο.
Η σταθερή παρουσία του \textit{Hour\_cos} προμηνύει τον ρόλο της ημερολογιακής δομής,
που κυριαρχεί στις πολυβηματικές στρατηγικές.

\begin{figure}[htbp]
  \centering
  \includegraphics[width=0.95\textwidth]{ch4_feature_importance/fig_fi_D_tf_load.png}
  \caption[Σημαντικότητα ανά μοντέλο: Teacher-Forcing, φορτίο]{Σημαντικότητα ανά
           μοντέλο: Teacher-Forcing, ηλεκτρικό φορτίο.}
  \label{fig:fi_D_tf_load}
\end{figure}

Στη Recursive (Σχ.~\ref{fig:fi_D_rec_load}) οι δύο αλγόριθμοι αποκλίνουν
έντονα. Το LightGBM περικόπτει το \textit{Load\_lag1} σε $7{,}3\%$,
αναδιανέμοντας στις κοντινές υστερήσεις (\textit{Load\_lag2}, \textit{Load\_lag6},
\textit{Load\_lag12}), όπως ακριβώς και στην τιμή. Το XGBoost με προγραμματισμένη
δειγματοληψία, αντιθέτως, διατηρεί το \textit{Load\_lag1} στο $84{,}4\%$: η ομαλότητα
του φορτίου κρατά αξιόπιστη ακόμη και την προβλεπόμενη προηγούμενη τιμή. Αυτό εξηγεί,
στο επίπεδο των χαρακτηριστικών, γιατί η αναδρομική στρατηγική πλησιάζει στατιστικά το
Teacher-Forcing για το φορτίο ($p=0{,}769$, §\ref{sec:dm_validation}), ενώ στην
ευμετάβλητη τιμή η ίδια αραίωση κοστίζει.

\begin{figure}[htbp]
  \centering
  \includegraphics[width=0.95\textwidth]{ch4_feature_importance/fig_fi_D_rec_load.png}
  \caption[Σημαντικότητα ανά μοντέλο: Recursive, φορτίο]{Σημαντικότητα ανά μοντέλο:
           Recursive, ηλεκτρικό φορτίο.}
  \label{fig:fi_D_rec_load}
\end{figure}

% ═════════════════════════════════════════════════════════════════════════════
\section{MIMO και Direct}
\label{sec:fi_mimo_direct}
% ═════════════════════════════════════════════════════════════════════════════

Στις δύο πολυβηματικές στρατηγικές η πρόβλεψη παράγεται χωρίς καμία αυτοπαλίνδρομη
ανατροφοδότηση: η μόνη πρόσφατη πληροφορία είναι οι παρατηρούμενες ενδοημερήσιες
υστερήσεις, οπότε οι εξωγενείς παράγοντες αποκτούν μεγαλύτερο σχετικό βάρος.

\begin{figure}[htbp]
  \centering
  \includegraphics[width=0.95\textwidth]{ch4_feature_importance/fig_fi_D_mimo_price.png}
  \caption[Σημαντικότητα ανά μοντέλο: MIMO, τιμή]{Σημαντικότητα ανά μοντέλο:
           MIMO ($H=24$), τιμή DAM.}
  \label{fig:fi_D_mimo_price}
\end{figure}

Χωρίς πρόσφατη τιμή, τα μοντέλα χωρίζονται σε δύο ομάδες. Τα δένδρα ενίσχυσης
διοχετεύουν τη βραχυπρόθεσμη πληροφορία στις παρατηρούμενες πυκνές υστερήσεις
4-23 (XGBoost: $53{,}5\%$ συλλογικά), που εδώ είναι καθαρή ιστορία και όχι
προβλέψεις, εξηγώντας τη βελτίωση $-7{,}9\%$ MAE έναντι του MIMO χωρίς πυκνές
υστερήσεις. Το Random Forest αντιθέτως αγκυρώνει στους κυλιόμενους μέσους 24h
($51{,}6\%$) και 168h ($16{,}6\%$), θυσιάζοντας οξύτητα στις αιχμές για ευστάθεια
επιπέδου, ενώ MLP και SVR στρέφονται στις ημερήσιες προβλέψεις παραγωγής και στον
καιρό (υγρασία, ακτινοβολία). Η ίδια απουσία αυτοπαλίνδρομου σήματος γεννά
αντίθετες στρατηγικές χαρακτηριστικών ανάλογα με τη μεροληψία του εκάστοτε
αλγορίθμου.

\begin{figure}[htbp]
  \centering
  \includegraphics[width=0.95\textwidth]{ch4_feature_importance/fig_fi_D_direct_price.png}
  \caption[Σημαντικότητα ανά μοντέλο: Direct, τιμή]{Σημαντικότητα ανά μοντέλο:
           Direct ($H=24$), τιμή DAM (LightGBM, XGBoost).}
  \label{fig:fi_D_direct_price}
\end{figure}

Το Direct αναπαράγει σχεδόν ταυτόσημα το MIMO: το XGBoost κρατά τις ενδοημερήσιες
υστερήσεις (\textit{Price\_lag6}, \textit{Price\_lag19}), το LightGBM στρέφεται
στις αιολικές υστερήσεις παραγωγής (48h, 24h) και στον κυλιόμενο μέσο 168h. Η
ταύτιση επιβεβαιώνει ότι MIMO και Direct διαφέρουν στη γεωμετρία εκπαίδευσης
(ένα κοινό έναντι ανεξάρτητων μοντέλων ανά ορίζοντα), όχι στην πληροφορία που
αξιοποιούν. Η επίπεδη κατανομή, με κανένα χαρακτηριστικό άνω του $5\%$,
αποτυπώνει την απουσία πρόσφατης τιμής: η πρόβλεψη στηρίζεται σε πολλά ασθενή αλλά
συμπληρωματικά σήματα.

Στο φορτίο, η απουσία πρόσφατης υστέρησης αναδεικνύει με προβλέψιμο τρόπο την
ημερολογιακή δομή. Στο MIMO (Σχ.~\ref{fig:fi_D_mimo_load}) το Random Forest αγκυρώνει
στον κυλιόμενο μέσο 24h (\textit{Load\_roll24}, $31{,}0\%$) και στα ωριαία ημερολογιακά
(\textit{Hour\_cos} $15{,}7\%$, \textit{Hour} $14{,}3\%$), ενώ το XGBoost διασπείρει το
βάρος στις παρατηρούμενες ενδοημερήσιες υστερήσεις (\textit{Load\_lag1/2/6/12}). Η
ντετερμινιστική ημερήσια και εβδομαδιαία περιοδικότητα του φορτίου υποκαθιστά την
πρόσφατη τιμή πολύ αποτελεσματικότερα απ' ό,τι στην τιμή DAM, όπου τον ρόλο αυτό
αναλαμβάνουν οι εξωγενείς παράγοντες (φυσικό αέριο, ΑΠΕ).

\begin{figure}[htbp]
  \centering
  \includegraphics[width=0.95\textwidth]{ch4_feature_importance/fig_fi_D_mimo_load.png}
  \caption[Σημαντικότητα ανά μοντέλο: MIMO, φορτίο]{Σημαντικότητα ανά μοντέλο:
           MIMO ($H=24$), ηλεκτρικό φορτίο.}
  \label{fig:fi_D_mimo_load}
\end{figure}

Το Direct (Σχ.~\ref{fig:fi_D_direct_load}) επαναλαμβάνει το μοτίβο: το XGBoost κρατά
τις ενδοημερήσιες υστερήσεις (\textit{Load\_lag1/6/12/3}), ενώ το LightGBM στρέφεται
στους κυλιόμενους μέσους (\textit{Load\_roll168}, \textit{Load\_roll24}), στην ημέρα
της εβδομάδας (\textit{Dow}, $3{,}3\%$) και στις αιολικές υστερήσεις
(\textit{Wind\_lag168}). Η ταύτιση MIMO-Direct, ίδια με την τιμή, επιβεβαιώνει ότι η
διαφορά τους είναι αμιγώς γεωμετρία εκπαίδευσης και όχι διαθέσιμη πληροφορία.

\begin{figure}[htbp]
  \centering
  \includegraphics[width=0.95\textwidth]{ch4_feature_importance/fig_fi_D_direct_load.png}
  \caption[Σημαντικότητα ανά μοντέλο: Direct, φορτίο]{Σημαντικότητα ανά μοντέλο:
           Direct ($H=24$), ηλεκτρικό φορτίο (LightGBM, XGBoost).}
  \label{fig:fi_D_direct_load}
\end{figure}

% ═════════════════════════════════════════════════════════════════════════════
\section{Συνολική εικόνα ανά στρατηγική}
\label{sec:fi_strategy}
% ═════════════════════════════════════════════════════════════════════════════

Συγκεντρώνοντας τους πρωταθλητές gain κάθε στρατηγικής σε ενιαία κλίμακα, η
μετατόπιση της βαρύτητας γίνεται άμεσα ορατή. Επειδή και οι τέσσερις στήλες
μετρούν την ίδια ποσότητα (gain), τα χρώματα είναι εδώ απευθείας συγκρίσιμα.

\begin{figure}[htbp]
  \centering
  \includegraphics[width=0.82\textwidth]{ch4_feature_importance/fig_fi_C_heatmap_price.png}
  \caption[Σημαντικότητα ανά στρατηγική, τιμή]{Σημαντικότητα δέκα κορυφαίων
           χαρακτηριστικών ανά στρατηγική (πρωταθλητές gain), τιμή DAM. Ενιαία
           κλίμακα \texttt{PowerNorm} ($\gamma=0{,}45$), ώστε να φαίνονται οι μικρές
           τιμές δίπλα στο κορυφαίο $88\%$.}
  \label{fig:fi_C_price}
\end{figure}

Το \textit{Price\_lag1} καταρρέει μονοτονικά, $88{,}3 \to 8{,}9 \to 4{,}9 \to
1{,}4\%$, καθώς η στρατηγική αφαιρεί σταδιακά την πρόσβαση στην πρόσφατη τιμή
(TF $\to$ Rec $\to$ MIMO $\to$ Direct). Το κενό που αφήνει το γεμίζουν διαφορετικά
σήματα σε κάθε στρατηγική: μακρινότερες υστερήσεις και κυλιόμενοι μέσοι στη
Recursive, ενδοημερήσιες υστερήσεις και αιολική παραγωγή στα πολυβηματικά σχήματα.
Η εικόνα δείχνει ότι η αυτοπαλίνδρομη κυριαρχία δεν είναι ιδιότητα της τιμής, αλλά
της διαθεσιμότητας πρόσφατης πληροφορίας που ορίζει η στρατηγική.

\begin{figure}[htbp]
  \centering
  \includegraphics[width=0.82\textwidth]{ch4_feature_importance/fig_fi_C_heatmap_load.png}
  \caption[Σημαντικότητα ανά στρατηγική, φορτίο]{Σημαντικότητα δέκα κορυφαίων
           χαρακτηριστικών ανά στρατηγική (πρωταθλητές gain), ηλεκτρικό φορτίο.
           Ίδια κλίμακα με το Σχ.~\ref{fig:fi_C_price}.}
  \label{fig:fi_C_load}
\end{figure}

Το φορτίο ακολουθεί την ίδια τροχιά κατάρρευσης, $92{,}4 \to 7{,}3 \to 6{,}9 \to
2{,}6\%$ για το \textit{Load\_lag1}, με μία ουσιώδη διαφορά στο τι αναδύεται.
Αντί για υστερήσεις παραγωγής, το βάρος μεταναστεύει στα ημερολογιακά
χαρακτηριστικά: η ώρα της ημέρας (\textit{Hour\_cos}) και η ημέρα της εβδομάδας
(\textit{Dow}) κρατούν σταθερό βάρος και κορυφώνονται στο Direct ($3{,}3\%$ η
\textit{Dow}). Σημειώνεται ότι η υψηλή συγκέντρωση στο TF
($92{,}4\%$, ελαφρώς πάνω από την τιμή) αφορά τους πρωταθλητές gain.η εποχική
δομή του φορτίου αναδεικνύεται στις μεθόδους που ευνοούν επίπεδες κατανομές (MDI,
permutation) και στις πολυβηματικές στρατηγικές.

% ═════════════════════════════════════════════════════════════════════════════
\section{Κατηγοριοποίηση χαρακτηριστικών}
\label{sec:fi_categories}
% ═════════════════════════════════════════════════════════════════════════════

Τα χαρακτηριστικά κατατάσσονται σε τρεις κατηγορίες, με βάρος που μετατοπίζεται
ανά στρατηγική.

\paragraph*{Α. Πρωτεύον αυτοπαλίνδρομο σήμα.}
Τα \textit{Price\_lag1} (ή \textit{Load\_lag1}) και \textit{Price\_lag24}
κυριαρχούν στο TF και στη Recursive, όπου το μοντέλο έχει πρόσβαση σε πρόσφατες
τιμές. Στο LightGBM-TF το \textit{Price\_lag1} εξηγεί μόνο του $88{,}3\%$ ενώ στη
Recursive πέφτει σε $8{,}9\%$ λόγω της αραίωσης.

\paragraph*{Β. Εξωγενείς παράγοντες αγοράς.}
Τα \textit{gas\_price}, \textit{co2\_price}, οι υστερήσεις αιολικής και ηλιακής
παραγωγής και οι ημερήσιες προβλέψεις παραγωγής αποκτούν μεγαλύτερο σχετικό βάρος
στο MIMO και στο Direct, όπου η αυτοπαλίνδρομη πληροφορία είναι περιορισμένη. Στο
LGBM-Direct οι αιολικές υστερήσεις παραγωγής ανέρχονται στις πρώτες θέσεις,
αντανακλώντας τη σύνδεση διείσδυσης ΑΠΕ και τιμής εκκαθάρισης.

\paragraph*{Γ. Ημερολογιακά και εποχικά.}
Τα \textit{Hour\_sin/cos}, \textit{Dow}, \textit{is\_holiday},
\textit{Price\_lag168} και \textit{Load\_lag168} έχουν σταθερή παρουσία στο φορτίο
σε όλες τις στρατηγικές και δευτερεύουσα στην τιμή. Εκφράζουν ντετερμινιστικές
δομές (εργάσιμες ώρες, αργίες, εβδομαδιαίος κύκλος), ισχυρότερες για το φορτίο.

Καθώς περνάμε από TF σε Recursive και σε MIMO, η εξάρτηση μετατοπίζεται από το
άμεσο \textit{Price\_lag1}, σε διασπορά προς μακρινότερες υστερήσεις και εξωγενείς
μεταβλητές, και τέλος σε εκτεταμένη χρήση των πυκνών υστερήσεων και των
ΑΠΕ/καυσίμων. Η συνέπεια είναι άμεση: η βέλτιστη επιλογή συνόλου χαρακτηριστικών
εξαρτάται από τη στρατηγική πρόβλεψης.

% ═════════════════════════════════════════════════════════════════════════════
\section{Σύνθεση και συμπεράσματα}
\label{sec:fi_conclusions}
% ═════════════════════════════════════════════════════════════════════════════

Η ανάλυση συμπυκνώνεται σε πέντε μηχανισμούς που εξηγούν τη συμπεριφορά κάθε
στρατηγικής.

\paragraph*{Teacher-Forcing: αυτοπαλίνδρομη κυριαρχία.}
Το \textit{Price\_lag1} συγκεντρώνει $88{,}3\%$ της σημαντικότητας στο LightGBM
της τιμής, το \textit{Load\_lag1} $92{,}4\%$ στο φορτίο. Η επιλογή είναι
στατιστικά αποδοτική αλλά εύθραυστη σε μεταβολή καθεστώτος, όπου η αυτοπαλίνδρομη
δυναμική αλλάζει ρυθμό, όπως φάνηκε στην κατάρρευση του SVR στο Q1 2026.

\paragraph*{Recursive: κατανεμημένο βάρος, εμφάνιση εξωγενών.}
Χωρίς πρόσβαση σε πραγματικές ενδοημερήσιες τιμές, το \textit{Price\_lag1}
υποχωρεί σε $8{,}9\%$ και το βάρος μεταφέρεται στα \textit{Price\_lag24} και
\textit{Price\_lag168}. Στο XGB-SS-Dense η προγραμματισμένη δειγματοληψία στρέφει
το μοντέλο στις πυκνές υστερήσεις 4-23 ($17{,}1\%$ συλλογικά).

\paragraph*{MIMO και Direct: στροφή στα εξωγενή.}
Χωρίς ενδιάμεσες πραγματικές τιμές, το XGB-MIMO Dense αξιοποιεί τις πυκνές
υστερήσεις σε ποσοστό $53{,}5\%$, ποσοστό που εξηγεί τη βελτίωση $-7{,}9\%$ MAE,
ενώ το RF-MIMO two-stage στηρίζεται στο \textit{Price\_roll24} ($51{,}6\%$) ως
άγκυρα. Το Direct αναπαράγει την ίδια δομή.

\paragraph*{Το παράδοξο των πυκνών υστερήσεων.}
Οι ίδιες πυκνές υστερήσεις βελτιώνουν το MIMO ($-22\%$ MLP, $-7{,}9\%$ XGB,
$-6{,}5\%$ LGBM) αλλά επιβαρύνουν τη Recursive ($+2{,}8\%$ LGBM). Στο MIMO οι
υστερήσεις είναι πάντα καθαρές ιστορικές τιμές. Στη Recursive γεμίζουν από
προβλέψεις, και ο θόρυβος διαδίδεται στα επόμενα βήματα. Η προγραμματισμένη
δειγματοληψία μετριάζει μερικώς το πρόβλημα (LGBM-SS-Dense: $-1{,}6\%$ έναντι
LGBM-SS).

\paragraph*{Φορτίο: ισχυρότερη εποχικότητα, ηπιότερη αραίωση.}
Υπό gain το φορτίο συγκεντρώνεται στο \textit{Load\_lag1} όσο και η τιμή
($92{,}4\%$ έναντι $88{,}3\%$), αλλά η εποχική του δομή το θωρακίζει στις
πολυβηματικές στρατηγικές: τα ημερολογιακά χαρακτηριστικά κρατούν σταθερό βάρος
ακόμη και χωρίς πρόσφατη υστέρηση. Η ομαλότερη δυναμική επιτρέπει στο Recursive
XGB-SS να ανταγωνίζεται στατιστικά το Teacher-Forcing ($p = 0{,}769$, DM test
Q1 2026 §\ref{sec:dm_validation}, Σχ.~\ref{fig:dm_load_wins_c6}), σε αντίθεση με την
τιμή, όπου η αραίωση κοστίζει.

Συνολικά, η σημαντικότητα είναι συνεπής εντός κάθε στρατηγικής, ανεξάρτητα από
τον αλγόριθμο. Η στρατηγική καθορίζει ποια πληροφορία είναι διαθέσιμη, και
δευτερευόντως ο αλγόριθμος καθορίζει πώς τη χρησιμοποιεί. Το ίδιο χαρακτηριστικό
($y_{t-1}$) περνά από $88\%$ στο TF σε $1$-$5\%$ στο MIMO: η επιλογή συνόλου
χαρακτηριστικών δεν είναι ανεξάρτητη από τον σχεδιασμό της πρόβλεψης.


